% Pandoc template for IEEE conference paper sections
\section{Related Work}

\textbf{Pairing-friendly curve taxonomy.} Freeman, Scott, and Teske
\cite{FreemanScottTeske2010} provide a comprehensive classification of
pairing-friendly curve constructions, including the Cocks-Pinch method,
MNT curves, and polynomial families (BN, BLS, KSS). Their work
establishes the theoretical framework but does not address
NTT-friendliness as a design criterion.

\textbf{BN curves and TNFS downgrade.} Barreto-Naehrig curves
\cite{BarretoNaehrig2005} achieve \(\rho = 1\) with embedding degree
12, making BN254 the standard for Ethereum's precompiles. However, Kim
and Barbulescu \cite{KimBarbulescu2016} demonstrated that the extended
Tower Number Field Sieve reduces BN254's effective security from 128-bit
to approximately 100-bit. This TNFS vulnerability motivates the move to
larger fields: our 1024-bit base field with \(k=18\) yields an
18,432-bit pairing target, providing substantial margin against all
known NFS variants.

\textbf{TNFS security updates.} Barbulescu and Duquesne
\cite{BarbulescuDuquesne2019} provided updated key size estimations
showing that many deployed pairing-friendly curves offer less security
than originally claimed. Their analysis confirms that embedding degree
18 with a 1024-bit prime provides \(\geq 250\)-bit security in the
pairing target field.

\textbf{BLS12-381.} Bowe \cite{BLS12_381} designed BLS12-381
specifically for zk-SNARKs, achieving two-adicity 32 in both fields
through careful parameter selection within the BLS12 family. Our work
achieves higher two-adicity (36) using the more flexible Cocks-Pinch
method, at the cost of a larger \(\rho\)-value (2.0 vs.~1.5).

\textbf{KSS-18 curves.} Kachisa, Schaefer, and Scott \cite{KSS2008}
construct pairing-friendly curves with embedding degree 18 using the
Brezing-Weng method \cite{BrezingWeng2005}, achieving
\(\rho \approx 1.33\). While KSS-18 offers better bandwidth efficiency
than Cocks-Pinch (\(\rho = 2\)), the parametric family does not
inherently guarantee NTT-friendliness, and the fixed polynomial
structure limits the ability to optimize for readability across a large
candidate space.

\textbf{Pasta curves.} The Pallas/Vesta cycle \cite{PastaCurves}
achieves NTT-friendliness through a 2-cycle of curves, enabling
recursive proof composition. Unlike our approach, Pasta curves are not
pairing-friendly and target a different application domain (Halo 2
recursive proofs).

\textbf{Nothing-up-my-sleeve numbers and SafeCurves.} The cryptographic
community has long advocated for transparent parameter generation to
demonstrate absence of backdoors. The SafeCurves project
\cite{SafeCurves} formalizes criteria including rigidity, twist
security, and completeness. Our readability scoring takes a
complementary approach: rather than deriving primes from a fixed public
constant, we select the most inspectable prime from a large space of
equally secure candidates, while the Cocks-Pinch seed provides full
reproducibility.

\textbf{Cocks-Pinch extensions.} Chen, Hong, and Wang
\cite{CocksPinchK5K8} extend the Cocks-Pinch method to embedding
degrees 5--8 with optimal Ate pairing computation. Our work focuses on
\(k=18\) with the novel addition of NTT constraints and readability
scoring.
